% ── sections/02_related_work.tex ──────────────────────────────────────────────

\section{Related Work}

\subsection{GNN-Based Cancer Driver Gene Prediction}

Identifying cancer driver genes computationally has been approached through
network-based and machine learning methods. Early work used network propagation
algorithms that diffused mutation signals across PPI networks, but these lacked
the capacity to integrate multiple omics data types. The introduction of graph
neural networks to this domain marked a significant advance.

EMOGI \cite{emogi2021} pioneered the integration of multi-omics node features
with PPI network topology using graph convolutional networks (GCNs),
demonstrating that GCN-based propagation over a single PPI graph substantially
outperforms feature-only baselines. The model encodes each gene as a node with
multi-omics features (mutation frequency, DNA methylation, gene expression, and
copy-number alteration) and learns to classify genes through neighbourhood
aggregation. EMGNN \cite{emgnn2023} extended this to a multilayer framework:
separate GNN encoders independently process each of multiple PPI networks, and a
meta-graph---whose nodes represent unique genes aggregated across networks---enables
cross-network evidence integration. Both methods employ Integrated Gradients for
feature-level interpretability, attributing predictions to individual omics
features.

Recent work has addressed specific limitations. SGCD \cite{sgcd2024} made the
important observation that PPI-based cancer driver prediction can be affected by
heterophily, with neighbours of cancer drivers often being non-cancer genes.
This finding motivates heterophily-aware architectures, as standard GNNs
implicitly assume homophily and may dilute cancer signals through neighbourhood
aggregation. DISHyper \cite{dishyper2024} incorporated Gene Ontology and disease
phenotype annotations as hypergraph structures, showing that functional knowledge
complements PPI topology. DGHNN \cite{dghnn2025} combined graph and hypergraph
encoders with a Feature Tokenizer Transformer and reported strong results on
EMOGI benchmark datasets. deepCDG \cite{deepcdg2025} introduced
weight-sharing across GCN layers with cross-omics attention for multi-omics
integration.

\subsection{Architectural Advances in Graph Learning}

Several general-purpose GNN innovations are relevant to our extensions. GPS
\cite{gps2022} combines local message passing with global multi-head
self-attention, overcoming the limited receptive field of standard GCNs.
GraphMAE \cite{graphmae2022} provides self-supervised pretraining via masked
feature reconstruction, exploiting unlabelled graph structure. Focal Loss
\cite{focal2017} and DropEdge \cite{dropedge2020} address class imbalance and
over-smoothing respectively---both common challenges in node classification.
PINNACLE \cite{pinnacle2024} provides pretrained protein embeddings from
multi-tissue PPI networks, while HIPGNN \cite{hipgnn2025} reframes anomaly
detection on graphs using spectral and spatial views. Hypergraph neural networks
\cite{hgnn2019} enable higher-order relationship modelling beyond pairwise edges,
relevant for pathway-level gene set encoding.

\subsection{Positioning of This Work}

Unlike prior single-contribution studies, our work provides a systematic pipeline
from reproduction through optimisation to interpretation. Key distinguishing
aspects include: (1) explicit gain decomposition separating data effects from
architectural effects; (2) systematic ablation of nine advanced GNN techniques
under a unified experimental protocol; (3) the finding that heterophily-aware
gating---rather than loss-function modifications or pretraining---provides the
largest single improvement. Table \ref{tab:relwork} positions our work against
the most relevant methods.

\begin{table}[ht]
  \centering
  \caption{Comparison of GNN-based cancer driver gene prediction methods.
           Multi-net: supports multiple PPI networks. Interp.: interpretability
           analysis. Hyper.: hypergraph encoding. AUPR values on CPDB test set
           where available; dashes indicate non-comparable protocols.}
  \label{tab:relwork}
  \small
  \begin{tabular}{lccccc}
    \toprule
    \textbf{Method} & \textbf{Year} & \textbf{Multi-net} & \textbf{Interp.} & \textbf{Hyper.} & \textbf{AUPR $\uparrow$} \\
    \midrule
    EMOGI \cite{emogi2021}      & 2021 & \texttimes & IG        & \texttimes & 0.73$^*$ \\
    EMGNN \cite{emgnn2023}      & 2023 & \checkmark & IG        & \texttimes & 0.809$^\dagger$ \\
    SGCD \cite{sgcd2024}        & 2025 & \texttimes & \texttimes & \texttimes & --- \\
    DISHyper \cite{dishyper2024} & 2024 & \texttimes & \texttimes & \checkmark & --- \\
    deepCDG \cite{deepcdg2025}  & 2025 & \texttimes & \texttimes & \texttimes & --- \\
    DGHNN \cite{dghnn2025}      & 2025 & \texttimes & \texttimes & \checkmark & --- \\
    \textbf{This work}          & 2026 & \checkmark & IG+GSEA   & \checkmark$^\ddagger$ & \textbf{0.824} \\
    \bottomrule
  \end{tabular}

  \vspace{0.3em}
  {\footnotesize $^*$Reported on CPDB. $^\dagger$Reported mean over 5 runs.
  $^\ddagger$Implemented, pending external data for evaluation.}
\end{table}
